% =====================================================================
\chapter{Cardinalidad, palomar y diagonal de Cantor}
\label{tema:6}
% =====================================================================

Las funciones inyectivas, sobreyectivas y biyectivas permiten comparar el tamaño de dos conjuntos sin contar sus elementos. En conjuntos finitos esta comparación da el principio del palomar; en conjuntos infinitos conduce a la distinción entre conjuntos numerables y no numerables y al argumento diagonal de Cantor. El tema sigue la sección~12.3 y el capítulo~14 de Hammack~\cite{hammack}, el capítulo~8 de Velleman~\cite{velleman}, las secciones~4.5, 7.1 y~14.8 de~\cite{mcs} y las transparencias «Cardinality» de Stanford CS103~\cite{cs103card}.

\section{Comparar tamaños con funciones}

\begin{definicion}[Equipotencia]
Dos conjuntos $A$ y $B$ tienen el \emph{mismo cardinal}, y se escribe $\abs{A}=\abs{B}$, si existe una biyección $A\to B$. Se escribe $\abs{A}\le\abs{B}$ si existe una función inyectiva $A\to B$.
\end{definicion}

Para conjuntos finitos, estas definiciones coinciden con contar elementos (\cref{prop:t3-finitos}). La relación «tener el mismo cardinal» es reflexiva ($\id_A$ es biyectiva), simétrica (la inversa de una biyección es biyectiva, \cref{prop:t4-inv-comp}) y transitiva (la composición de biyecciones es biyectiva, \cref{teo:t4-comp}).

\addtocounter{definicion}{1}%

\section{El principio del palomar}

\begin{teorema}[Principio del palomar]\label{teo:t6-palomar}
Sean $A$ y $B$ finitos y $f\colon A\to B$.
\begin{enumerate}
  \item Si $\abs{A}>\abs{B}$, $f$ no es inyectiva: existen $a\neq a'$ con $f(a)=f(a')$.
  \item (Forma general.) Si $\abs{A}>k\abs{B}$ con $k\in\N$, alguna fibra tiene al menos $k+1$ elementos. En particular, alguna fibra tiene al menos $\lceil\abs{A}/\abs{B}\rceil$ elementos.
\end{enumerate}
\end{teorema}

El nombre procede de la imagen de $\abs{A}$ palomas que se refugian en $\abs{B}$ palomares: si hay más palomas que palomares, en algún palomar hay al menos dos. La dificultad en las aplicaciones está en elegir bien las «palomas» y los «palomares».

\addtocounter{definicion}{1}%

\begin{proposicion}[Límite de la compresión sin pérdida]\label{prop:t6-compresion}
Para $n\ge1$ no existe ninguna función inyectiva de $\{0,1\}^{n}$ en el conjunto $\{0,1\}^{<n}=\bigcup_{k=0}^{n-1}\{0,1\}^{k}$ de cadenas binarias de longitud menor que $n$.
\end{proposicion}

Un compresor sin pérdida es una función inyectiva (debe poder descomprimirse). La \cref{prop:t6-compresion} dice que ningún compresor acorta \emph{todas} las cadenas de longitud $n$: si acorta algunas, alarga o deja igual otras. El módulo~6 cuantifica esta idea con la entropía.

\section{Conjuntos numerables}

\begin{definicion}[Numerable]
Un conjunto $A$ es \emph{infinito numerable} si $\abs{A}=\abs{\N}$, es decir, si sus elementos pueden listarse como $a_0,a_1,a_2,\dots$ sin repeticiones y sin omitir ninguno. Es \emph{numerable} si es finito o infinito numerable, y \emph{no numerable} en otro caso.
\end{definicion}

\begin{proposicion}\label{prop:t6-numerables}
Los conjuntos $\Z$, $\N\times\N$ y $\Sigma^{*}$ (con $\Sigma$ finito y no vacío) son infinitos numerables.
\end{proposicion}

\addtocounter{figure}{1}%

\begin{proposicion}\label{prop:t6-criterios}
Sea $A$ un conjunto no vacío. Son equivalentes: (i) $A$ es numerable; (ii) existe una función inyectiva $A\to\N$; (iii) existe una función sobreyectiva $\N\to A$. Además, todo subconjunto de un conjunto numerable es numerable, y la unión de una familia numerable de conjuntos numerables es numerable.
\end{proposicion}

Se admite este resultado sin demostración completa (véase~\cite[§14.2]{hammack} y~\cite[§8.2]{velleman}); algunas de sus partes se piden en los \cref{eje:6-subconjN,eje:6-producto}. Como consecuencia, $\Q$ es numerable (\cref{eje:6-Q}).

\section{El argumento diagonal de Cantor}

\begin{teorema}[Cantor]\label{teo:t6-cantor}
Para todo conjunto $A$ no existe ninguna función sobreyectiva $A\to\Pa(A)$. En consecuencia, $\abs{A}\neq\abs{\Pa(A)}$.
\end{teorema}

\begin{corolario}\label{cor:t6-nonum}
$\Pa(\N)$ y el conjunto $\{0,1\}^{\N}$ de sucesiones infinitas de ceros y unos no son numerables.
\end{corolario}

El nombre «diagonal» se entiende con sucesiones: dada una lista $s_0,s_1,s_2,\dots$ de sucesiones binarias, la sucesión $d$ con $d_k=1-s_k(k)$ difiere de $s_k$ en la posición $k$, luego no está en la lista. El mismo argumento, con desarrollos decimales y evitando las colas de nueves, prueba que $\R$ no es numerable~\cite[§14.2]{hammack}.

\addtocounter{definicion}{1}%

\addtocounter{definicion}{1}%

